\section{Introduction}

William of Malmesbury, writing his \textit{Gesta Regum Anglorum} in about 1125, records
that Eilmer, an older monk of his own abbey, had in his youth fixed wings to his hands and
feet, launched from a tower, and travelled ``more than a furlong'' before he fell, broke both
legs, and was lame for the rest of his life \autocite{williamMalmesbury_gesta}.
Eilmer, William adds, used to say that he fell because he had forgotten to fit a tail.
Lynn White's study of the passage established it as the earliest credible account of a
gliding attempt in Western Europe \autocite{white1961eilmer}.

The short account names a launch point, a distance, a wind, a failure, an injury, and a
cause, and undergraduate physics can check each one.
Could Eilmer's flight have happened at Malmesbury, and is a missing tail enough to explain
the fall?
No source describes Eilmer's apparatus beyond its wings, so the models sample its unknown
parameters.

As a teaching case for undergraduate mechanics, the problem assumes calculus, Newtonian
mechanics, and elementary probability, and suits students from the second year onwards and
their instructors.
Answering it needs the energy balance of a glider at launch, flight in a moving air mass,
lift over a slope, the neutral-point criterion for pitch stability, the growth rate of an
unstable mode, Monte Carlo estimation with confidence intervals, and the mechanics of a
feet-first impact.
Each step uses public data and a few lines of code, gives a number to check against the
text, and becomes an exercise in \cref{sec:teaching}.
